\documentclass[12pt,a4paper]{article}

% ------------------ PACKAGES ------------------
\usepackage[utf8]{inputenc}
\usepackage[T1]{fontenc}
\usepackage{geometry}
\usepackage{setspace}
\usepackage{graphicx}
\usepackage{longtable}
\usepackage{array}
\usepackage{booktabs}
\usepackage{hyperref}
\usepackage{listings}
\usepackage{xcolor}
\usepackage{enumitem}
\usepackage{fancyhdr}
\usepackage{listings}
\usepackage{xcolor}

\lstdefinelanguage{JavaScript}{
  keywords={const, let, var, function, return, if, else, for, while, break, continue, new},
  keywordstyle=\color{blue}\bfseries,
  ndkeywords={class, export, boolean, throw, implements, import, this},
  ndkeywordstyle=\color{purple}\bfseries,
  identifierstyle=\color{black},
  sensitive=true,
  comment=[l]{//},
  morecomment=[s]{/*}{*/},
  commentstyle=\color{gray}\ttfamily,
  stringstyle=\color{red}\ttfamily,
  morestring=[b]',
  morestring=[b]"
}


% ------------------ PAGE SETUP ------------------
\geometry{margin=1in}
\onehalfspacing
\pagestyle{fancy}
\fancyhf{}
\fancyhead[L]{Kigali Shaolin Temple Management System}
\fancyhead[R]{System Blueprint}
\fancyfoot[C]{\thepage}

% ------------------ CODE STYLE ------------------
\definecolor{codegray}{gray}{0.95}
\lstset{
  backgroundcolor=\color{codegray},
  basicstyle=\ttfamily\small,
  frame=single,
  breaklines=true
}

% ------------------ DOCUMENT ------------------
\begin{document}

% ------------------ TITLE PAGE ------------------
\begin{titlepage}
    \centering
    \vspace*{2cm}
    {\Huge \textbf{Kigali Shaolin Temple Management System}}\\[0.5cm]
    {\Large Frontend Website, Admin Portal (CMS), and Backend API}\\[0.3cm]
    {\large \textbf{Version 1.1.0}}\\[1.5cm]
    
    \textbf{Purpose:}\\
    A unified professional blueprint combining system analysis, UI/UX design,
    CMS control, and backend logic.\\[1.5cm]
    
    \vfill
    {\large \textbf{Prepared for:} Technical Documentation}\\
    {\large \textbf{Year:} 2025}
\end{titlepage}

\tableofcontents
\newpage

% ------------------------------------------------
\section{General Project Overview}

\textbf{Project Name:} Kigali Shaolin Temple Management System

\textbf{Purpose:}  
To promote Shaolin culture in Rwanda, manage clubs, programs, blogs, and allow administrators
to dynamically control website content through a secure Content Management System (CMS).

This document serves as a master reference for:
\begin{itemize}
    \item System development
    \item Technical documentation
    \item Project proposal or defense
    \item Team onboarding
\end{itemize}

% ------------------------------------------------
\section{System Architecture}

The system follows a modern, scalable, production-ready architecture.

\begin{verbatim}
[ Frontend (React + Vite) ]
            |
         REST API
            |
[ Backend (Node.js + Express) ]
            |
[ Admin Portal (CMS Dashboard) ]
            |
[ Database (PostgreSQL + Sequelize) ]
\end{verbatim}

\subsection*{Architecture Description}
\begin{itemize}
    \item Frontend: Public-facing website for visitors
    \item Admin Portal: Secure CMS for content management
    \item Backend API: Business logic, authentication, and data access
    \item Database: Persistent storage
\end{itemize}

% ------------------------------------------------
\section{Branding and UI Design System (Tailwind)}

\subsection{Color Palette}

\begin{center}
\begin{tabular}{ll}
\toprule
Role & Color \\
\midrule
Primary & \#7B1E1E (Shaolin Red) \\
Dark & \#111111 \\
Accent & \#C9A24D (Heritage Gold) \\
Text & \#FFFFFF \\
Muted & \#9CA3AF \\
\bottomrule
\end{tabular}
\end{center}

\subsection{Tailwind Configuration}
\begin{lstlisting}[language=JavaScript]
extend: {
  colors: {
    primary: '#7B1E1E',
    dark: '#0A0A0A',
    gold: '#C4A747',
  }
}
\end{lstlisting}

\subsection{Design Aesthetics (Premium Matrix)}
The system employs a high-end visual language referred to as the "Premium Matrix" aesthetic:
\begin{itemize}
    \item \textbf{Glassmorphism}: Interactive cards use \texttt{bg-white/70 backdrop-blur-xl} for a modern, layered feel.
    \item \textbf{Cinematic Overlays}: Hero sections and media displays use deep gradients and subtle glows.
    \item \textbf{Premium Shadows}: Custom shadows (\texttt{shadow-premium}) provide depth without clutter.
    \item \textbf{Interactive Micro-animations}: Smooth hover effects and transitions are applied globally.
\end{itemize}

\subsection{Typography}

\begin{center}
\begin{tabular}{ll}
\toprule
Usage & Font \\
\midrule
Headings & Playfair Display (Serif) \\
Body Text & Inter / Poppins (Sans-serif) \\
\bottomrule
\end{tabular}
\end{center}

% ------------------------------------------------
\section{Frontend Website (Public)}

\subsection{Design Goals}
\begin{itemize}
    \item Cinematic visuals
    \item Martial arts strength and discipline
    \item Traditional yet modern appearance
    \item Fully CMS-driven content
    \item Mobile-first responsiveness
\end{itemize}

% ------------------------------------------------
\section{Frontend Pages and Sections}

\subsection{Header and Navigation}
\begin{itemize}
    \item Logo
    \item Navigation: Home, About, Blog, Clubs, Contact
    \item Language switcher (RW / EN / FR)
\end{itemize}

CMS Control: None (static navigation)

\subsection{Hero Section}
\begin{itemize}
    \item Full-screen background image
    \item Dark overlay
    \item Title, subtitle, CTA button
    \item Carousel slider
\end{itemize}

CMS Controlled:
\begin{itemize}
    \item Title and subtitle
    \item Background images
    \item Button text and link
\end{itemize}

\subsection{About Us Section}
\begin{itemize}
    \item Image and text layout
    \item Founder quotation
\end{itemize}

CMS Controlled:
\begin{itemize}
    \item About text
    \item Founder name
    \item Images
    \item Year founded
\end{itemize}

\subsection{Programs and Training}
\begin{itemize}
    \item Online Training
    \item In-Person Training
    \item Global Retreats
\end{itemize}

Each program contains:
\begin{itemize}
    \item Image
    \item Title
    \item Description
    \item Status (active or inactive)
\end{itemize}

\subsection{Blogs}
\begin{itemize}
    \item Featured blog
    \item Blog listing
    \item Blog detail page
\end{itemize}

\subsection{Clubs}
\begin{itemize}
    \item Club name
    \item Location
    \item Instructor
    \item Description
    \item Images
\end{itemize}

\subsection{Contact Us}
\begin{itemize}
    \item Name
    \item Email
    \item Message
\end{itemize}

\subsection{Gallery}
\begin{itemize}
    \item Photo and video media
    \item Categorization by event
    \item Date and title metadata
\end{itemize}

Backend Requirements:
\begin{itemize}
    \item Message storage
    \item Admin replies
    \item Email notifications
\end{itemize}

% ------------------------------------------------
\section{Admin Portal (CMS Dashboard)}

\subsection{Authentication and Roles}

\begin{center}
\begin{tabular}{ll}
\toprule
Role & Description \\
\midrule
Super Admin & Complete system control and audit access \\
Content Manager & Global content management across all modules \\
Blogger & Focused access: creation and management of own blogs \\
Editor & Standard content editing and review \\
\bottomrule
\end{tabular}
\end{center}

\subsection{Dashboard Modules}
\begin{itemize}
    \item Dashboard overview and system statistics
    \item Homepage (Hero Slides) manager
    \item About manager
    \item Programs manager
    \item Blog manager
    \item Clubs and Tutorials manager
    \item Gallery (Media) manager
    \item Messages manager
    \item Audit Logs (Activity tracking)
    \item Language manager
    \item Users and roles
\end{itemize}

% ------------------------------------------------
\section{Backend (API and Logic)}

\subsection{Core Responsibilities}
\begin{itemize}
    \item Authentication and authorization using JWT
    \item Content APIs
    \item Image uploads
    \item Email services
    \item Role-based access control
\end{itemize}

\subsection{Technology Stack}

Backend:
\begin{itemize}
    \item Node.js with Express
    \item JWT authentication and Authorization middleware
    \item Multer and Cloudinary for media storage
    \item Nodemailer for email invitations and password resets
    \item Swagger (OpenAPI) for interactive API documentation
    \item Audit Logging for administrative actions
    \item \textbf{API Resilience}: Centralized \texttt{apiClient.js} implementation
\end{itemize}

\subsection{API Interceptors and Security}
The frontend communication layer is hardened with:
\begin{itemize}
    \item \textbf{Request Interceptors}: Automatic injection of JWT Bearer tokens from local storage.
    \item \textbf{Response Interceptors}: Global handling of 401 Unauthorized errors with automatic cleanup and redirect to login.
    \item \textbf{Centralized Configuration}: Environment-based Base URL management for seamless dev/prod switching.
\end{itemize}
\subsection{Media Upload Architecture}
The system supports two upload modes controlled by the \texttt{UPLOAD\_MODE} environment variable:

\begin{center}
\begin{tabular}{lll}
\toprule
Mode & Storage & Persistence \\
\midrule
\texttt{local} & Server filesystem (\texttt{/uploads}) & Ephemeral (lost on restart) \\
\texttt{cloud} & Cloudinary CDN & Permanent \\
\bottomrule
\end{tabular}
\end{center}

\textbf{Production Requirement}: Always use \texttt{UPLOAD\_MODE=cloud} with valid Cloudinary credentials on Render (and any platform with ephemeral storage), as local files are wiped on container restart.

Cloudinary is configured once at module load in \texttt{uploadUtils.js}:
\begin{lstlisting}[language=JavaScript]
cloudinary.config({
  cloud_name: process.env.CLOUD_NAME,
  api_key: process.env.API_KEY,
  api_secret: process.env.API_SECRET,
});
\end{lstlisting}


% ------------------------------------------------
\section{Database Models}

\begin{lstlisting}[language=SQL]
User(id, name, email, password, role, image, reset_token, created_at)

HeroSlide(id, title, subtitle, image, button_text, button_link)

About(id, content, founder_name, image)

Program(id, name, description, image, status)

Blog(id, title, content, image, category, published, created_at)

Club(id, name, description, location, image)

ClubTutorial(id, club_id, title, video_url)

Gallery(id, title, description, url, media_type, event_date, category)

ContactMessage(id, name, email, message, read, created_at)

AuditLog(id, user_id, action, entity, entity_id, details, ip_address)
\end{lstlisting}

% ------------------------------------------------
\section{Advanced System Features}

\subsection{Audit Logging}
All administrative actions (Create, Update, Delete, Restore) are logged in the system. Each log entry captures:
\begin{itemize}
    \item The administrator who performed the action
    \item The type of action and the entity affected
    \item A timestamp and the administrator's IP address
\end{itemize}

\subsection{Data Recovery and Soft Deletes}
The system implements soft deletes (paranoid mode), allowing administrators to restore accidentally deleted content. Permanent deletion is only possible through a deliberate "Force Delete" action.

\subsection{Interactive API Documentation}
The backend includes a Swagger (OpenAPI) interface, providing real-time documentation and a testing environment for all API endpoints.

\subsection{ID Obfuscation (UUID Encryption)}
To enhance security and prevent exposure of internal database structure, the system implements ID obfuscation using Base64URL encoding:

\begin{itemize}
    \item All public-facing UUIDs are encoded using Base64URL (URL-safe encryption)
    \item Backend transparently decodes obfuscated IDs via middleware
    \item Prevents enumeration attacks and information disclosure
    \item Maintains backward compatibility with existing UUID structure
\end{itemize}

Implementation:
\begin{lstlisting}[language=JavaScript]
// Encoding: UUID -> Base64URL
const encodeUUID = (uuid) => {
  const hex = uuid.replace(/-/g, '');
  const buffer = Buffer.from(hex, 'hex');
  return buffer.toString('base64')
    .replace(/\+/g, '-')
    .replace(/\//g, '_')
    .replace(/=/g, '');
};

// Decoding middleware applied to all routes
router.get('/:id', decodeParam('id'), controller.getById);
\end{lstlisting}

\subsection{Database Schema Standardization}
All Sequelize models have been standardized with explicit $snake_cas$e timestamp mappings to ensure database consistency:

\begin{itemize}
    \item Explicit mapping: \texttt{createdAt: '$created_at$', updatedAt: '$updated_at$'}
    \item Resolves Sequelize ordering and query inconsistencies
    \item Prevents 500 Internal Server Errors from column name mismatches
    \item Applied to all models: Club, Program, Gallery, About, HeroSlide, User, Blog, ContactMessage, AuditLog
\end{itemize}

\subsection{Frontend Stability Enhancements}
The admin portal includes robust error handling and defensive programming practices:

\begin{itemize}
    \item Optional chaining for safe property access (\texttt{pagination?.totalItems})
    \item Default values for undefined data states
    \item Pagination protection against crashes during network latency
    \item Error boundaries for graceful degradation
\end{itemize}

% ------------------------------------------------
\section{Docker Infrastructure and Deployment}

The system is fully containerized using Docker, enabling consistent deployment and simplified environment management across development, testing, and production.

\subsection{Container Services}
The ecosystem consists of three orchestrated services. Production images are published on Docker Hub under the \texttt{andremugabo} registry with the \texttt{:prod} tag:
\begin{itemize}
    \item \textbf{kst-backend}: \texttt{andremugabo/kst-backend:prod} --- Node.js API with Cloudinary integration and custom entrypoint for DB synchronization and seeding.
    \item \textbf{kst-frontend}: \texttt{andremugabo/kst-frontend:prod} --- Production-optimized public website served via Nginx.
    \item \textbf{kst-admin}: \texttt{andremugabo/kst-admin:prod} --- Production-optimized CMS portal served via Nginx.
\end{itemize}

\subsection{Deployment Workflow}
The entire stack can be initialized with a single command:
\begin{lstlisting}[language=bash]
docker compose up -d --build
\end{lstlisting}

\subsection{Production Build Commands}
For production environments, the Frontend and Admin apps are built with the live API URL \textbf{and upload URL} baked in at build time (Vite resolves \texttt{import.meta.env} at build time, not runtime):
\begin{lstlisting}[language=bash]
cd kst.rw
# Build and push all images (uses push-images.sh)
bash push-images.sh

# Manual example for Frontend:
docker build \
  --build-arg VITE_API_URL=https://kst-backend.onrender.com/api \
  --build-arg VITE_UPLOAD_URL=https://kst-backend.onrender.com \
  -t andremugabo/kst-frontend:prod ./kst-frontend
docker push andremugabo/kst-frontend:prod
\end{lstlisting}

\subsection{SPA Routing and Nginx}

Both the Admin Portal and Frontend are configured with a custom Nginx configuration to support Single Page Application (SPA) routing, ensuring that all deep-linked URLs correctly resolve to \texttt{index.html}.

\subsection{Automated Initialization}
The backend container utilizes a dedicated \texttt{docker-entrypoint.sh} script that:
\begin{enumerate}
    \item Waits for the database service to be healthy.
    \item Runs \texttt{npm run db:reset} to synchronize models and seed initial administrative users.
    \item Starts the production Express server.
\end{enumerate}

% ------------------------------------------------
\section{Production Hosting (Render.com)}

For production environments, the system is designed to be hosted on Render.com using their Infrastructure-as-Code (Blueprint) feature.

\subsection{Render Blueprint (render.yaml)}
A \texttt{render.yaml} file is provided, which provisions all three web services using Docker images.
Key environment variables for the backend:

\begin{center}
\begin{tabular}{ll}
\toprule
Variable & Purpose \\
\midrule
\texttt{NODE\_ENV} & Set to \texttt{production} \\
\texttt{DATABASE\_URL} & Neon.tech PostgreSQL connection string \\
\texttt{BACKEND\_URL} & Public backend URL for file path construction \\
\texttt{UPLOAD\_MODE} & \texttt{cloud} (required for persistence) \\
\texttt{CLOUD\_NAME} & Cloudinary cloud name \\
\texttt{API\_KEY} & Cloudinary API key \\
\texttt{API\_SECRET} & Cloudinary API secret \\
\texttt{JWT\_SECRET} & Auto-generated by Render \\
\bottomrule
\end{tabular}
\end{center}

\noindent Key environment variables for the frontend and admin:

\begin{center}
\begin{tabular}{ll}
\toprule
Variable & Purpose \\
\midrule
\texttt{VITE\_API\_URL} & Backend API base URL \\
\texttt{VITE\_UPLOAD\_URL} & Backend base URL for media resolution \\
\bottomrule
\end{tabular}
\end{center}

\textbf{Note}: \texttt{VITE\_UPLOAD\_URL} must also be passed as a Docker \texttt{--build-arg} since Vite bakes environment variables at build time.

\subsection{Production Data Safety}
The system implements environment-aware startup logic. When \texttt{NODE\_ENV=production} is detected:
\begin{itemize}
    \item Full database resets (\texttt{db:reset}) are disabled.
    \item Safe synchronization (\texttt{db:sync}) is performed using Sequelize's \texttt{alter: true} mode.
    \item The seeder (\texttt{seed.js}) runs in upsert mode --- it creates records if absent, and updates stale placeholder URLs if they already exist.
\end{itemize}

% ------------------------------------------------
\section{Resolved Technical Issues (v1.1.0)}

\begin{center}
\begin{tabular}{p{5cm}p{9cm}}
\toprule
Issue & Resolution \\
\midrule
\texttt{ERR\_NAME\_NOT\_RESOLVED} for placeholder images & Replaced \texttt{via.placeholder.com} with \texttt{placehold.co} in seed data and all components \\
Images disappear on page refresh & Root cause: \texttt{UPLOAD\_MODE=cloud} with placeholder Cloudinary credentials silently fell back to ephemeral local storage. Fixed by calling \texttt{cloudinary.config()} at module load and providing real credentials. \\
Backend fails to start on Render & Removed hardcoded \texttt{PORT: 3005} from \texttt{render.yaml}; Render assigns ports dynamically. \\
All API requests blocked by CORS & Service name mismatch between Render hostnames and \texttt{allowedOrigins} in \texttt{server.js}. Fixed by correcting the CORS list. \\
Malformed image URLs & Fixed leading/trailing slash handling in \texttt{getMediaPath} in \texttt{apiClient.js}. \\
No image URLs in production frontend & \texttt{VITE\_UPLOAD\_URL} was missing as a Docker \texttt{build-arg}; added to both Dockerfiles and \texttt{push-images.sh}. \\
\bottomrule
\end{tabular}
\end{center}

% ------------------------------------------------
\section{Non-Functional Requirements}

\begin{itemize}
    \item Mobile-first responsive design
    \item Search engine optimization
    \item Secure admin routes
    \item Optimized images
    \item Fast loading performance
\end{itemize}

% ------------------------------------------------
\section{Final Summary}

The Kigali Shaolin Temple Management System consists of three integrated layers:

\begin{center}
\begin{tabular}{lll}
\toprule
Layer & Role & Live URL \\
\midrule
Frontend & Public Website & \href{https://kst-frontend.onrender.com}{kst-frontend.onrender.com} \\
Admin Portal & CMS Dashboard & \href{https://kst-admin.onrender.com}{kst-admin.onrender.com} \\
Backend & API Engine & \href{https://kst-backend.onrender.com}{kst-backend.onrender.com} \\
\bottomrule
\end{tabular}
\end{center}

This design:
\begin{itemize}
    \item Covers the majority of UI requirements from the provided design
    \item Uses Tailwind CSS for a modern and professional interface
    \item Enables scalability through CMS control
\end{itemize}

\end{document}
